% Eigenständiger Prosabeweis; alle nummerierten Aussagen bleiben in Band 45.
\section*{Eine Rechenregel verrät, was enthalten ist}
\hypertarget{semilattice.reading.example}{}

Eine Rechentafel sagt uns, welches Ergebnis zu zwei Eingaben gehört.
Eine Ordnung sagt uns, welches Element unter welchem anderen liegt.
Auf den ersten Blick sind das verschiedene Arten von Information.
Bei bestimmten Rechenregeln lässt sich jedoch die eine vollständig
aus der anderen gewinnen. Ein kleines Mengenbeispiel zeigt schon
den entscheidenden Gedanken.

Seien \(a\) und \(b\) verschieden und \(S=\{a,b\}\). Die
Potenzmenge von \(S\), also die Menge aller ihrer Teilmengen, ist
\[
  A=\powerset(S)=\{\varnothing,\{a\},\{b\},S\}.
\]
Unsere Eingaben sind die vier \emph{Teilmengen}, nicht die beiden
Elemente \(a,b\). Als Rechenoperation verwenden wir die Vereinigung.
Zum Beispiel ergeben \(\{a\}\) und \(\{b\}\) zusammen die Menge
\(S\); die Vereinigung von \(\{a\}\) mit \(S\) ergibt wieder
\(S\). Alle Ergebnisse stehen in der folgenden Tafel.

\SemilatticeExampleDiagram

Der Zusammenhang zwischen beiden Darstellungen lautet
\begin{equation}
  X\subseteq Y\quad\Longleftrightarrow\quad X\cup Y=Y.
  \tag{V}\label{semilattice-reading-union}
\end{equation}
Ist jedes Element von \(X\) bereits in \(Y\), fügt die
Vereinigung nichts hinzu; deshalb ist \(X\cup Y=Y\). Gilt
umgekehrt diese Gleichheit, so liegt jedes Element von \(X\)
in \(X\cup Y\) und damit in \(Y\). Also ist \(X\subseteq Y\).
Das beweist beide Richtungen von
\FormulaRefAuto{SubsetIffUnionEqualsRight}[thm] in Band~45.

Man kann die Ordnung somit an der Tafel ablesen: \(X\) liegt
unter \(Y\), wenn in der Zeile \(X\), Spalte \(Y\) wieder
der Spaltenkopf \(Y\) steht. Beispielsweise gilt
\(\{a\}\subseteq S\), weil \(\{a\}\cup S=S\). Dagegen
ist \(\{a\}\cup\{b\}=S\) weder \(\{a\}\) noch \(\{b\}\).
Keine der beiden Einermengen ist in der anderen enthalten.

Auch die Rechentafel lässt sich aus dem Bild zurückgewinnen. Um
\(X\cup Y\) zu finden, suchen wir das kleinste Element, das
oberhalb von \(X\) und von \(Y\) liegt. Für die beiden
Einermengen ist das \(S\), für \(\varnothing\) und
\(\{a\}\) ist es \(\{a\}\). Was hier an vier Elementen leicht
zu sehen ist, wird im Folgenden für eine beliebige Menge bewiesen.
Die vier Tabellen zum Vereinigungsbeispiel stehen in den
\SemilatticeProofsLink{}; die Begründung des Zusammenhangs bleibt
vollständig in diesem Lesetext.

\par\Needspace{13\baselineskip}
\section*{Welche Rechengesetze werden gebraucht?}
\hypertarget{semilattice.reading.laws}{}

Wir beginnen mit einer Menge \(A\) und einer binären Operation
\[
  \vee\colon A\times A\longrightarrow A.
\]
Zu jedem geordneten Paar \((x,y)\) gehört damit genau ein Wert
\(x\vee y\), und dieser Wert liegt wieder in \(A\).
Die letzte Forderung heißt \emph{Abgeschlossenheit}. Sie erlaubt,
ein Zwischenergebnis erneut mit einem Element von \(A\) zu
verknüpfen. Das Zeichen \(\vee\) steht hier für eine Operation
auf \(A\), nicht für die logische Verknüpfung zweier Aussagen.

Zusätzlich verlangen wir für alle \(x,y,z\in A\) drei Gesetze:
\[
  \begin{aligned}
    (x\vee y)\vee z&=x\vee(y\vee z)
      &&\text{Assoziativität},\\
    x\vee y&=y\vee x
      &&\text{Kommutativität},\\
    x\vee x&=x
      &&\text{Idempotenz}.
  \end{aligned}
\]
Die erste Gleichung erlaubt das Umklammern von drei Eingaben.
Die zweite erlaubt das Vertauschen der Eingaben. Die dritte sagt,
dass die Verknüpfung eines Elements mit sich selbst nichts ändert.
Eine solche Struktur heißt \emph{Halbverband} oder
\emph{Halbgitter}. Sie ist eine kommutative idempotente Halbgruppe;
die Definition steht unter \FormulaRefAuto{Semilattice}[def].
Ein neutrales Element gehört nicht zu diesen Voraussetzungen.

Die Vereinigung erfüllt alle Forderungen. Sind \(X,Y\subseteq S\),
so enthält \(X\cup Y\) weiterhin nur Elemente von \(S\).
Das ist die Abgeschlossenheit aus
\FormulaRefAuto{PowerSetUnionClosure}[thm]. Ein Element liegt
genau dann in \((X\cup Y)\cup Z\), wenn es in mindestens
einer der Mengen \(X,Y,Z\) liegt; dieselbe Bedingung beschreibt
\(X\cup(Y\cup Z)\). Das beweist die Assoziativität.
Die Bedingungen \glqq in \(X\) oder in \(Y\)\grqq{} und
\glqq in \(Y\) oder in \(X\)\grqq{} sind gleichwertig,
also ist die Vereinigung kommutativ. Schließlich ist
\glqq in \(X\) oder in \(X\)\grqq{} gleichbedeutend mit
\glqq in \(X\)\grqq{}, woraus \(X\cup X=X\) folgt.
Damit ist auch \FormulaRefAuto{PowerSetUnionSemilattice}[thm]
erklärt. Die Argumente gelten für jede Menge \(S\).

Nach dem Vorbild von \eqref{semilattice-reading-union} definieren
wir nun eine Relation auf \(A\):
\begin{equation}
  x\leq_\vee y\quad\Longleftrightarrow\quad x\vee y=y
  \qquad(x,y\in A).
  \tag{O}\label{semilattice-reading-order}
\end{equation}
Der Index erinnert daran, dass diese Relation erst aus \(\vee\)
gewonnen wird. In Band~45 heißt sie der \emph{supremal induzierte
Relationsoperator}, siehe \FormulaRefAuto{JoinOrder}[def].
Dort wird die Abhängigkeit von der Operation im Kontext mitgeführt.
Hier hilft der Index später beim Vergleich mit einer bereits
gegebenen Ordnung.

Als Menge geordneter Paare ist unsere Relation
\[
  R_\vee=\{(x,y)\in A\times A\mid x\vee y=y\}.
\]
Aus \(x\leq_\vee y\) folgen deshalb sowohl \(x,y\in A\)
als auch \(x\vee y=y\). Das sind die drei vorbereitenden
Aussagen \FormulaRefAuto{JoinOrderFirstCarrier}[thm],
\FormulaRefAuto{JoinOrderSecondCarrier}[thm] und
\FormulaRefAuto{JoinOrderEvaluation}[thm]. Die Beschränkung auf
\(A\) sichert bei jeder Anwendung, dass die Operation definiert ist.

\medskip\noindent\textbf{Das Ziel.}
Wir werden beweisen, dass \(\leq_\vee\) eine partielle Ordnung
ist und \(x\vee y\) jeweils die kleinste gemeinsame obere
Schranke von \(x\) und \(y\) liefert. Anschließend kehren wir
den Weg um: Eine partielle Ordnung, in der jedes Elementpaar
eine solche Schranke besitzt, bestimmt einen Halbverband.
Beide Übergänge gewinnen die ursprünglichen Daten zurück.

\par\Needspace{10\baselineskip}
\section*{Aus den Rechengesetzen wird eine Ordnung}
\hypertarget{semilattice.reading.order}{}

Eine partielle Ordnung auf \(A\) muss reflexiv, antisymmetrisch
und transitiv sein. Wir prüfen diese Eigenschaften direkt anhand
von \eqref{semilattice-reading-order}. Alle im Folgenden genannten
Elemente liegen in \(A\).

\medskip\noindent\textbf{Jedes Element liegt unter sich selbst.}
Die Idempotenz liefert \(x\vee x=x\). Nach der Definition
der Relation bedeutet das \(x\leq_\vee x\). Damit ist die
Reflexivität bewiesen, also
\FormulaRefAuto{JoinOrderReflexive}[thm].

\medskip\noindent\textbf{Gegenseitige Vergleiche erzwingen Gleichheit.}
Es gelte \(x\leq_\vee y\) und \(y\leq_\vee x\).
Übersetzt in die Operation heißt das
\[
  x\vee y=y,\qquad y\vee x=x.
\]
Wegen der Kommutativität stimmen die linken Seiten überein.
Also erhalten wir die Gleichheitskette
\[
  x=y\vee x=x\vee y=y.
\]
Das ist die Antisymmetrie aus
\FormulaRefAuto{JoinOrderAntisymmetric}[thm]. Sie verbietet
nicht, dass ein Element mit sich selbst vergleichbar ist.
Sie sagt, dass zwei \emph{verschiedene} Elemente nicht zugleich
in beiden Richtungen untereinander liegen können.

\medskip\noindent\textbf{Vergleiche lassen sich fortsetzen.}
Es gelte nun \(x\leq_\vee y\) und \(y\leq_\vee z\),
also \(x\vee y=y\) und \(y\vee z=z\). Dann gilt
\[
  \begin{aligned}
    x\vee z
      &=x\vee(y\vee z)\\
      &=(x\vee y)\vee z\\
      &=y\vee z\\
      &=z.
  \end{aligned}
\]
In der ersten Zeile setzen wir \(z=y\vee z\) ein. Die zweite
verwendet die Assoziativität; die dritte setzt \(x\vee y=y\)
ein. Die letzte ist die zweite Voraussetzung. Nach
\eqref{semilattice-reading-order} folgt \(x\leq_\vee z\).
Damit ist die Transitivität aus
\FormulaRefAuto{JoinOrderTransitive}[thm] bewiesen.

Die drei Prüfungen zeigen gemeinsam
\FormulaRefAuto{JoinSemilatticeInducesOrder}[thm]:
\((A,\leq_\vee)\) ist eine partiell geordnete Menge.
Dabei hatten die drei Rechengesetze jeweils eine sichtbare Aufgabe:
Idempotenz liefert Reflexivität, Kommutativität ermöglicht die
Antisymmetrie, und Assoziativität verbindet zwei Vergleiche zu
einem dritten. Der formale Abschluss steht unter
\SemilatticeProofLink{JoinSemilatticeInducesOrder}.

Eine partielle Ordnung muss nicht \emph{total} sein. Totalität
würde zusätzlich verlangen, dass für jedes \(x,y\in A\)
mindestens einer der Vergleiche \(x\leq y\), \(y\leq x\)
gilt. Die beiden Einermengen unseres Beispiels zeigen, dass diese
Forderung hier nicht folgt. Da \(a\neq b\), ist
\(\{a\}\cup\{b\}=S\) von beiden Einermengen verschieden.
Weder \(\{a\}\leq_\cup\{b\}\) noch
\(\{b\}\leq_\cup\{a\}\) gilt. Sie sind
\emph{unvergleichbar}, obwohl ihre Vereinigung eindeutig bestimmt ist.

Für sämtliche Teilmengen von \(S\) gilt dagegen nach
\eqref{semilattice-reading-union}
\[
  X\leq_\cup Y\quad\Longleftrightarrow\quad X\subseteq Y.
\]
Die aus der Vereinigung gewonnene Ordnung ist also genau die
Inklusionsordnung, wie
\FormulaRefAuto{PowerSetUnionOrderIsInclusion}[thm] festhält.
Unsere allgemeine Konstruktion hat in diesem Beispiel keine neue
Art des Vergleichs erzeugt, sondern die vertraute wiedergefunden.

\par\Needspace{11\baselineskip}
\section*{Warum das Ergebnis die kleinste obere Schranke ist}
\hypertarget{semilattice.reading.supremum}{}

Ein Element \(u\in A\) heißt \emph{gemeinsame obere Schranke}
von \(x\) und \(y\), wenn \(x\leq_\vee u\) und
\(y\leq_\vee u\) gelten. Ein \emph{Supremum} des Paares
ist eine solche obere Schranke \(s\), die unter jeder anderen
gemeinsamen oberen Schranke liegt. Es müssen also zwei verschiedene
Forderungen erfüllt sein:
\[
  x\leq_\vee s,\quad y\leq_\vee s,
  \qquad
  \bigl(x\leq_\vee u\ \land\ y\leq_\vee u\bigr)
       \Longrightarrow s\leq_\vee u.
\]
Die zweite Forderung gilt für \emph{jedes} \(u\in A\).
Sie macht aus irgendeiner oberen Schranke die kleinste.

Wir setzen \(s=x\vee y\). Wegen der Abgeschlossenheit liegt
\(s\) in \(A\). Zunächst zeigen wir, dass beide Eingaben unter
dem Ergebnis liegen. Für die erste gilt
\[
  x\vee s=x\vee(x\vee y)=(x\vee x)\vee y=x\vee y=s.
\]
Hier benutzen wir Assoziativität und Idempotenz. Nach der Definition
von \(\leq_\vee\) ist damit \(x\leq_\vee s\) bewiesen.
Für die zweite Eingabe gilt entsprechend
\[
  y\vee s
    =y\vee(x\vee y)
    =y\vee(y\vee x)
    =(y\vee y)\vee x
    =y\vee x
    =s.
\]
Zusätzlich wurde die Kommutativität verwendet. Also ist auch
\(y\leq_\vee s\). Das sind
\FormulaRefAuto{JoinUpperLeft}[thm] und
\FormulaRefAuto{JoinUpperRight}[thm].

Nun sei \(u\) eine beliebige gemeinsame obere Schranke. Dann
gelten \(x\vee u=u\) und \(y\vee u=u\). Daraus folgt
\[
  s\vee u=(x\vee y)\vee u=x\vee(y\vee u)=x\vee u=u.
\]
Somit ist \(s\leq_\vee u\). Das beweist
\FormulaRefAuto{JoinLeastUpper}[thm] und schließt den Nachweis ab:
\begin{equation}
  x\vee y=\sup_{\leq_\vee}\{x,y\}.
  \tag{S}\label{semilattice-reading-supremum}
\end{equation}
Der Index am Supremumszeichen nennt die betrachtete Ordnung;
der zugrunde liegende Träger ist weiterhin \(A\).
Die Aussage steht in Band~45 als
\FormulaRefAuto{JoinIsPairSupremum}[thm], ihre tabellarische
Ableitung unter \SemilatticeProofLink{JoinIsPairSupremum}.

\medskip\noindent\textbf{Supremum, Minimum und Maximum.}
Das Supremum ist das \emph{Minimum der Menge aller oberen
Schranken}. Es ist im Allgemeinen weder ein Minimum noch ein
Maximum der ursprünglichen Zweiermenge. Im Beispiel betrachten wir
die Zweiermenge
\[
  B=\bigl\{\{a\},\{b\}\bigr\}\subseteq\powerset(S).
\]
Ihr Supremum in \((\powerset(S),\subseteq)\) ist \(S\).
Die Menge \(S\) ist jedoch kein Element von \(B\): Die beiden
Elemente von \(B\) sind die Einermengen. Deshalb hat \(B\)
kein Maximum. Es hat auch kein Minimum, denn keine der beiden
Einermengen liegt unter der anderen. Dass \(\varnothing\) unter
beiden liegt, macht es nicht zum Minimum von \(B\); auch
\(\varnothing\) gehört nicht zu \(B\).

Ist hingegen \(x\leq_\vee y\), dann liegt \(y\) über
beiden Elementen und gehört selbst zu \(\{x,y\}\). Es ist
dann deren Maximum und zugleich ihr Supremum. Genau diesen Fall
erkennt die Gleichung \(x\vee y=y\).

Auch \emph{kleinste} und \emph{minimale} obere Schranke dürfen
nicht verwechselt werden. Bei einer minimalen oberen Schranke gibt
es keine echt kleinere obere Schranke unterhalb von ihr. Andere
obere Schranken könnten aber unvergleichbar mit ihr sein.
Die kleinste obere Schranke liegt dagegen unter \emph{allen}
oberen Schranken. Gerade diese stärkere Eigenschaft haben wir
mit dem beliebigen \(u\) bewiesen.

\par\Needspace{12\baselineskip}
\section*{Die Rückrichtung: Aus einer Ordnung wird eine Operation}
\hypertarget{semilattice.reading.reverse}{}

Nun beginnen wir von der anderen Seite. Gegeben seien eine Menge
\(A\) und eine partielle Ordnung \(\leq\) auf \(A\).
Wir setzen voraus, dass für jedes \(x,y\in A\) ein Supremum
von \(\{x,y\}\) in \(A\) existiert. Eine Operation ist in
diesem Schritt noch nicht gegeben. Wir wollen sie aus den Suprema
gewinnen.

Zuerst ist die Eindeutigkeit zu klären. Sind \(s\) und \(t\)
beide Suprema desselben Paares, ist \(t\) eine obere Schranke.
Die kleinste-Schranke-Eigenschaft von \(s\) liefert daher
\(s\leq t\). Ebenso ist \(s\) eine obere Schranke, sodass
\(t\leq s\) folgt. Wegen der Antisymmetrie gilt \(s=t\).
Die Voraussetzung liefert also zu jedem Paar genau ein Supremum.

Deshalb dürfen wir definieren
\begin{equation}
  x\vee y:=\sup_{\leq}\{x,y\}
  \qquad(x,y\in A).
  \tag{P}\label{semilattice-reading-pair-operation}
\end{equation}
Dies legt tatsächlich eine binäre Funktion auf \(A\) fest.
Ihr Graph ist die Menge aller \(((x,y),m)\in(A\times A)\times A\),
für die \(m\) die beiden Eigenschaften
\[
  x\leq m,\quad y\leq m,
  \qquad
  \forall u\in A\;
    \bigl((x\leq u\land y\leq u)\Longrightarrow m\leq u\bigr)
\]
erfüllt. Diese Menge ist eine durch die angegebenen Bedingungen
bestimmte Teilmenge von \((A\times A)\times A\).
Existenz der Suprema bedeutet: Über jedem Argument \((x,y)\)
liegt mindestens ein Wert. Der Eindeutigkeitsbeweis bedeutet:
Über diesem Argument liegt höchstens ein Wert. Zusammen ergeben
beide Aussagen eine Funktion \(A\times A\to A\).
Ein Auswahlaxiom wird nicht benötigt, denn bei keinem Paar muss
zwischen mehreren möglichen Werten gewählt werden.

Dies ist die durch Paarsuprema bestimmte Operation aus
\FormulaRefAuto{PairJoinOperation}[def]. Ihre Werte liegen in
\(A\), weil ein Supremum in der gegebenen Ordnung ein Element
von \(A\) sein muss. Abgeschlossenheit ist damit bereits gesichert.
Die drei Rechengesetze sind noch zu beweisen.

Für diesen Beweis halten wir fest, was die Definition unmittelbar
liefert:
\begin{equation}
  \begin{gathered}
    x\leq x\vee y,\qquad y\leq x\vee y,\\
    (x\leq u\ \land\ y\leq u)
      \quad\Longrightarrow\quad x\vee y\leq u.
  \end{gathered}
  \tag{U}\label{semilattice-reading-upper-rules}
\end{equation}
Dies sind \FormulaRefAuto{PairJoinUpperLeft}[thm],
\FormulaRefAuto{PairJoinUpperRight}[thm] und
\FormulaRefAuto{PairJoinLeastUpper}[thm]. Anders als in der
Hinrichtung wurden diese Regeln nicht aus Rechengesetzen gewonnen:
Sie sind jetzt der Inhalt der Definition von \(x\vee y\).
Sämtliche Vergleiche in \eqref{semilattice-reading-upper-rules}
verwenden die \emph{ursprünglich gegebene} Ordnung \(\leq\).

Im Hasse-Diagramm unseres Beispiels existiert für jedes Paar ein
kleinstes gemeinsames Element darüber. Beginnen wir nur mit dieser
Inklusionsordnung, weist \eqref{semilattice-reading-pair-operation}
den beiden Einermengen den Wert \(S\) zu. Dem Paar
\((\varnothing,\{a\})\) weist sie \(\{a\}\) zu.
So entstehen die ersten Einträge der ursprünglichen Tafel.

\par\Needspace{13\baselineskip}
\section*{Weshalb diese Operation die drei Rechengesetze erfüllt}
\hypertarget{semilattice.reading.associativity}{}

Wir verwenden weiterhin nur die partielle Ordnung \(\leq\)
und die drei Schrankenregeln
\eqref{semilattice-reading-upper-rules}. Insbesondere dürfen
wir beim Rechnen mit der neu definierten Operation noch nicht
umklammern.

\medskip\noindent\textbf{Idempotenz.}
Das Element \(x\) ist eine obere Schranke von \(\{x,x\}\),
denn \(x\leq x\). Jede obere Schranke dieses Paares liegt
definitionsgemäß über \(x\). Also ist \(x\) selbst das
Supremum. Die Definition der Operation liefert
\[
  x\vee x=x.
\]
Das ist \FormulaRefAuto{PairJoinIdempotent}[thm].

\medskip\noindent\textbf{Kommutativität.}
Die Mengen \(\{x,y\}\) und \(\{y,x\}\) sind gleich.
Ein Element liegt genau dann über \(x\) und \(y\), wenn
es über \(y\) und \(x\) liegt. Beide Paare besitzen daher
dieselbe Menge oberer Schranken und dasselbe eindeutige Supremum:
\[
  x\vee y=y\vee x.
\]
Das beweist \FormulaRefAuto{PairJoinCommutative}[thm].

\medskip\noindent\textbf{Assoziativität.}
Hier müssen wir zwei zunächst verschiedene Rechnungen vergleichen.
Für feste \(x,y,z\in A\) setzen wir
\[
  p:=(x\vee y)\vee z,\qquad q:=x\vee(y\vee z).
\]
Die beiden Werte existieren und liegen in \(A\), weil die
Operation abgeschlossen ist. Wir zeigen, dass jeder von ihnen
die kleinste gemeinsame obere Schranke der drei Elemente
\(x,y,z\) ist. Ein Supremum für drei Elemente wird dabei
nicht zusätzlich vorausgesetzt; wir weisen seine Eigenschaften nach.

Zuerst betrachten wir \(p\). Die Schrankenregeln geben
\[
  x\leq x\vee y,\quad y\leq x\vee y,\qquad
  x\vee y\leq p,\quad z\leq p.
\]
Aus der Transitivität der gegebenen Ordnung folgen
\(x\leq p\) und \(y\leq p\). Zusammen mit \(z\leq p\)
ist \(p\) also eine gemeinsame obere Schranke aller drei Elemente.

Sei nun \(u\in A\) irgendeine gemeinsame obere Schranke von
\(x,y,z\). Aus \(x\leq u\) und \(y\leq u\) folgt
zuerst \(x\vee y\leq u\). Zusammen mit \(z\leq u\)
liefert dieselbe Schrankenregel im nächsten Schritt
\[
  p=(x\vee y)\vee z\leq u.
\]
Damit liegt \(p\) unter jeder gemeinsamen oberen Schranke
von \(x,y,z\).

Für \(q\) führen wir beide Schritte ebenfalls aus.
Es gelten
\[
  y\leq y\vee z,\quad z\leq y\vee z,\qquad
  x\leq q,\quad y\vee z\leq q.
\]
Transitivität liefert \(y\leq q\) und \(z\leq q\).
Mit \(x\leq q\) ist also auch \(q\) eine gemeinsame
obere Schranke der drei Elemente. Ist \(u\) wiederum eine
beliebige gemeinsame obere Schranke, so folgt aus
\(y\leq u\), \(z\leq u\) zunächst \(y\vee z\leq u\).
Zusammen mit \(x\leq u\) erhalten wir
\[
  q=x\vee(y\vee z)\leq u.
\]
Auch \(q\) liegt unter jeder gemeinsamen oberen Schranke.

Da \(q\) selbst eine solche Schranke ist, folgt aus der für
\(p\) bewiesenen Eigenschaft \(p\leq q\). Da \(p\)
ebenfalls eine solche Schranke ist, folgt umgekehrt \(q\leq p\).
Antisymmetrie liefert \(p=q\), also
\[
  (x\vee y)\vee z=x\vee(y\vee z).
\]
Das ist \FormulaRefAuto{PairJoinAssociative}[thm]. In diesem
Beweis wurde keine Klammer verschoben. Beide Klammerungen wurden
getrennt durch ihre Ordnungseigenschaft charakterisiert und erst
am Schluss miteinander identifiziert. Die formale Ableitung steht
unter \SemilatticeProofLink{PairJoinAssociative}.

Am Mengenbeispiel lässt sich der Gedanke mit
\(X=\{a\}\), \(Y=\{b\}\), \(Z=\varnothing\) verfolgen.
Links wird zuerst das Supremum der Einermengen gebildet, also
\(S\), und anschließend das Supremum von \(S\) und
\(\varnothing\), wiederum \(S\). Rechts ergibt das innere
Paar zunächst \(\{b\}\), das äußere dann \(S\).
Beide Ergebnisse sind die kleinste Teilmenge von \(S\), die
alle drei Eingabemengen enthält.

Abgeschlossenheit und alle drei Rechengesetze sind nun bewiesen.
Die durch Paarsuprema definierte Operation macht \(A\) zu einem
Halbverband. Das ist
\FormulaRefAuto{PairJoinCharacterization}[thm]; der formale
Zusammenschluss steht unter
\SemilatticeProofLink{PairJoinCharacterization}.

\par\Needspace{12\baselineskip}
\section*{Beide Beschreibungen gewinnen einander zurück}
\hypertarget{semilattice.reading.recovery}{}

Wir haben zwei Übergänge konstruiert. Um eine vollständige
Entsprechung zu erhalten, müssen wir noch zeigen, dass ein
Hin- und Rückweg nichts verändert. Dies betrifft einmal die
Operation und einmal die Ordnung.

\medskip\noindent\textbf{Von der Operation zur Ordnung und zurück.}
Wir starten mit einem Halbverband \((A,\vee)\) und bilden
die Ordnung \(\leq_\vee\) durch
\eqref{semilattice-reading-order}. Aus deren Paarsuprema
definieren wir eine neue Operation, die wir vorübergehend
\(\diamond\) nennen:
\[
  x\diamond y:=\sup_{\leq_\vee}\{x,y\}.
\]
Nach \eqref{semilattice-reading-supremum} ist dieses Supremum
gerade \(x\vee y\). Also gilt für jedes \((x,y)\in A\times A\)
\[
  x\diamond y=x\vee y.
\]
Beide Funktionen haben denselben Definitionsbereich und stimmen
an jedem Argument überein. Sie sind deshalb dieselbe Operation.
Die ursprüngliche Rechentafel ist vollständig zurückgewonnen.

\medskip\noindent\textbf{Von der Ordnung zur Operation und zurück.}
Nun starten wir mit einer partiellen Ordnung \(\leq\), in der
alle Paarsuprema existieren. Daraus bilden wir \(\vee\) gemäß
\eqref{semilattice-reading-pair-operation}. Weil diese Operation
ein Halbverband ist, definiert
\[
  x\leq_\vee y\quad\Longleftrightarrow\quad x\vee y=y
\]
eine partielle Ordnung. Noch tragen die alte Ordnung \(\leq\)
und die neu gewonnene Ordnung \(\leq_\vee\) verschiedene Zeichen.
Für ihre Gleichheit zeigen wir beide Vergleichsrichtungen.

Es gelte zuerst \(x\leq y\) in der alten Ordnung.
Wegen \(y\leq y\) ist \(y\) eine gemeinsame obere Schranke
von \(x\) und \(y\). Die Supremumseigenschaft liefert
\(x\vee y\leq y\); zugleich gilt stets \(y\leq x\vee y\).
Antisymmetrie ergibt \(x\vee y=y\), also \(x\leq_\vee y\).

Es gelte umgekehrt \(x\leq_\vee y\). Dann ist
\(x\vee y=y\). Die erste Schrankenregel in der ursprünglichen
Ordnung besagt \(x\leq x\vee y\). Einsetzen der Gleichheit
ergibt \(x\leq y\). Damit ist
\[
  x\leq y\quad\Longleftrightarrow\quad x\leq_\vee y
  \qquad(x,y\in A)
\]
bewiesen. Die beiden Relationen sind dieselbe Teilmenge von
\(A\times A\). Dies ist
\FormulaRefAuto{PairJoinRecoversOrder}[thm]; der formale Nachweis steht unter
\SemilatticeProofLink{PairJoinRecoversOrder}.

\SemilatticeCorrespondenceDiagram

In unserem Eingangsbeispiel ist die Verbindung besonders konkret.
Aus der Vereinigungstafel liest man durch \(X\cup Y=Y\)
die Inklusionsordnung ab. Aus der Inklusionsordnung gewinnt man
die Tafel zurück, denn \(X\cup Y\) enthält \(X\) und
\(Y\); jede Teilmenge \(U\subseteq S\), die beide enthält,
enthält auch \(X\cup Y\). Also ist die Vereinigung genau
das geforderte Supremum. Tafel und Hasse-Diagramm beschreiben
dieselbe Struktur mit verschiedenen Mitteln.

\medskip\noindent\textbf{Wie weit reicht das Ergebnis?}
Die Menge \(A\) braucht nicht endlich zu sein. Ebenso wenig
muss die Ordnung total sein, wie schon die beiden unvergleichbaren
Einermengen zeigen. Erforderlich sind genau die Paarsuprema.
Ein kleinstes Element folgt daraus nicht: Auf \(\mathbb Z\)
mit der üblichen Ordnung ist
\[
  m\vee n=\max\{m,n\}
\]
stets definiert. Nach unserem Satz ist dies eine Halbverbandsoperation.
Ein kleinstes ganzzahliges Element gibt es dennoch nicht, denn
zu jedem \(n\) liegt \(n-1\) echt darunter.

Ein kleinstes Element \(0\), falls es vorhanden ist, wäre für
\(\vee\) neutral: Aus \(0\leq x\) folgt \(0\vee x=x\),
und wegen der Kommutativität auch \mbox{\(x\vee0=x\)}. Umgekehrt
wäre jedes neutrale Element \(e\) ein kleinstes Element der
induzierten Ordnung, weil \(e\vee x=x\) gerade
\(e\leq_\vee x\) bedeutet. Eine solche Null ist daher eine
zusätzliche Eigenschaft, keine versteckte Voraussetzung des Satzes.
Im Mengenbeispiel übernimmt \(\varnothing\) diese Rolle.

Aus Paarsuprema entstehen durch wiederholtes Verknüpfen auch die
Suprema aller endlichen nichtleeren Mengen: Nach dem ersten
Element fügt man jeweils eines hinzu; die Schrankenregeln zeigen
bei jedem Schritt, dass der neue Wert das Supremum der bisher
berücksichtigten Elemente ist. Für die leere Menge bräuchte man
dagegen ein kleinstes Element. Denn jedes Element ist obere
Schranke der leeren Menge; ihr Supremum müsste deshalb unter
allen Elementen von \(A\) liegen.

Auch für unendliche Teilmengen garantiert der Satz kein Supremum.
Die Menge aller ganzen Zahlen hat in \(\mathbb Z\) nicht einmal
eine obere Schranke, obwohl jedes Paar ein Maximum besitzt.
Die kleine Vereinigungstafel lässt diese Grenze nicht erkennen:
Eine Potenzmenge hat zusätzliche Eigenschaften, die für allgemeine
Halbverbände nicht vorausgesetzt werden dürfen.

Schließlich kann man die Richtung einer Ordnung umkehren.
Die kleinste obere Schranke wird dann zur größten unteren Schranke,
dem \emph{Infimum}. Dieselbe Halbverbandsoperation lässt sich
in der dualen Ordnung infimal lesen. Auf einer Potenzmenge liefert
außerdem der Durchschnitt das Infimum bezüglich der Inklusion.
Die gemeinsame Theorie von Suprema und Infima führt zu den
Verbänden; der Einstieg dazu steht unter \SemilatticeMainLink{}.
Die 23 formalen Ableitungen dieses Lesebogens sind in den
\SemilatticeProofsLink{} zusammengeführt, während ihre Aussagen
und Kennungen in Band~45 erhalten bleiben.

\clearpage
\section*{Wie die Verbindung historisch entstand}
\hypertarget{semilattice.reading.history}{}

Die Verbindung von Rechenregeln und Ordnung hat mehrere historische
Wurzeln. Ein Zugang führte über die Algebra der Logik. Ernst Schröder
behandelte 1890 im ersten Band seiner \emph{Vorlesungen über die
Algebra der Logik} das Rechnen mit Klassen und die Beziehungen
zwischen ihnen.\footnote{Ernst Schröder,
\href{https://www.deutschestextarchiv.de/book/show/schroeder_logik01_1890}{\emph{Vorlesungen über die Algebra der Logik}},
Bd.~1, Leipzig 1890. Dedekind verweist in der folgenden Quelle,
S.~112, ausdrücklich auf Schröders Behandlung der entsprechenden Gesetze.}
In heutiger Mengensprache begegnen hier Vereinigung, Durchschnitt
und Enthaltensein. Unser Eingangsbeispiel knüpft daran an:
Ob eine Klasse schon in einer anderen enthalten ist, lässt sich
daran erkennen, dass ihre Vereinigung die andere nicht vergrößert.

Ein weiterer Zugang kam aus Richard Dedekinds Zahlentheorie.
An positiven ganzen Zahlen lässt sich die zugrunde liegende
Idee unmittelbar sehen. Ordnet man sie nach Teilbarkeit, so
übernimmt das kleinste gemeinsame Vielfache die Rolle des
Supremums, der größte gemeinsame Teiler die des Infimums.
In unserer Schreibweise lautet der Zusammenhang
\[
  a\mid b\quad\Longleftrightarrow\quad
  \operatorname{kgV}(a,b)=b.
\]
So ist etwa \(\operatorname{kgV}(6,10)=30\): Beide Zahlen
teilen \(30\), und \(30\) teilt jedes gemeinsame Vielfache.
Das ist dieselbe Schrankenidee, die wir an der Vereinigung
von Mengen untersucht haben.

Dedekind löste solche gemeinsamen Gesetze von den einzelnen
Beispielen. In seiner Arbeit von 1897 über Zerlegungen von Zahlen
führt er \emph{Dualgruppen} als Systeme beliebiger Dinge mit
zwei Verknüpfungen und festgelegten Gesetzen ein. Er nennt unter
anderem Moduln, Ideale und Untergruppen als Beispiele und
vergleicht die Gesetze ausdrücklich mit denen der Logik.
Dabei berichtet er, dass ihn die Theorie seiner Moduln zu
diesen Fragen geführt habe. Im Aufsatz von 1900 über drei
Moduln führt er aus den Verknüpfungen eine verallgemeinerte
Teilbarkeit ein und stellt ihre Reflexivität, Antisymmetrie
und Transitivität heraus.\footnote{Richard Dedekind,
\href{https://rcin.org.pl/Content/140714/PDF/WA35_171681_15883-2_Art9.pdf}{\emph{Über Zerlegungen von Zahlen durch ihre größten gemeinsamen Teiler}}
(1897), nachgedruckt in \emph{Gesammelte mathematische Werke},
Bd.~II, S.~103--147; insbesondere \S\,3--4, S.~108--114;
sowie \href{https://rcin.org.pl/Content/140716/PDF/WA35_171699_15883-2_Art11.pdf}{\emph{Über die von drei Moduln erzeugte Dualgruppe}}
(1900), ebenda, S.~236--271, insbesondere \S\,1, S.~236--238.}
Die Abstraktion hält also fest, was verschiedene Arten von
Objekten gemeinsam haben: Man kann sie verknüpfen und daraus
ablesen, welches unter welchem liegt.

\par\Needspace{12\baselineskip}
Garrett Birkhoff untersuchte solche Strukturen in den 1930er
Jahren systematisch als Mittel der abstrakten Algebra. Sein
Aufsatz von 1933 über die Verbindung von Unteralgebren fragt
nach Folgerungen und Anwendungen gemeinsamer Axiome. In einer
Arbeit von 1935 definiert er eine abstrakte Inklusionsbeziehung
ausdrücklich durch eine Gleichung der Form
\mbox{\(a\cup b=b\)}.\footnote{Garrett Birkhoff,
\href{https://doi.org/10.1017/S0305004100011464}{\emph{On the combination of subalgebras}},
\emph{Proceedings of the Cambridge Philosophical Society}~29
(1933), S.~441--464; sowie
\href{https://doi.org/10.1017/S0305004100013463}{\emph{On the structure of abstract algebras}},
ebenda~31 (1935), S.~433--454, insbesondere S.~435, Fußnote~\S.}
Hier erscheint der Leitgedanke dieser Lesefassung unmittelbar:
Eine Verknüpfung bestimmt eine Ordnung.

Diese Entwicklung betrifft vor allem die Theorie der Verbände
mit zwei Verknüpfungen. Unsere Lesefassung betrachtet gezielt
die Seite der Suprema. Sie zeigt, dass schon eine einzige
assoziative, kommutative und idempotente Operation die
entsprechende Ordnung vollständig festlegt. Der Beweisgang
ordnet die Gedanken nach ihren logischen Abhängigkeiten;
die historischen Beispiele erklären, weshalb man nach einer
solchen gemeinsamen Struktur gesucht hat.

\par\Needspace{16\baselineskip}
\section*{Schlussbemerkung}
\hypertarget{semilattice.reading.conclusion}{}

Ausgangspunkt war eine kleine Vereinigungstafel. Ihr entscheidender
Zug ist, dass man das Enthaltensein an einer Rechnung erkennen
kann. Die drei Rechengesetze tragen diesen Gedanken auf eine
beliebige Menge weiter:
\[
  x\leq_\vee y\quad\Longleftrightarrow\quad x\vee y=y.
\]
Die so entstehende Ordnung hat für jedes Paar ein Supremum,
und dieses ist gerade das Ergebnis der ursprünglichen Operation.
Umgekehrt liefern die Paarsuprema einer Ordnung eine
Halbverbandsoperation. Beide Übergänge gewinnen ihre jeweilige
Ausgangsstruktur exakt zurück.

Damit sind Rechentafel und Ordnungsbild zwei vollständige
Beschreibungen derselben Struktur. Man kann eine algebraische
Gleichung als Aussage über kleinste obere Schranken lesen und
eine solche Ordnungseigenschaft wieder in eine Rechenregel
übersetzen. Endlichkeit oder Totalität sind dafür nicht nötig.
Ein kleinstes Element und Suprema beliebiger unendlicher Mengen
sind zusätzliche Eigenschaften. Der Gewinn liegt gerade darin,
die gemeinsame Struktur genau zu erkennen und ihre Grenzen
ebenso genau festzuhalten.
